%% ma: L11-B17-0067
%% lop: 11
%% chuong: 5
%% bai: 17
%% dang: 11.17.5
%% loai: DS
%% muc: [TH, TH, VD, VD]
%% dap_an: "SSĐĐ"
%% nguon: "tai-lieu-hoc-tap-mon-toan-lop-11-theo-chuong-trinh-giao-duc-pho-thong-2018.pdf (Lớp Toán Thầy Duy, Phạm Lê Duy, THPT Chu Văn An), Chương 3 Bài 3, Phần 2, Câu 7 (tr.417)"
%% kien_thuc: "Lớp 11 Bài 17 (hàm số liên tục)"
%% trang_thai: cho-duyet
%% ly_do: "Dạng mới đề xuất (11.17.x), chờ giáo viên duyệt dạng; tài liệu gốc không có lời giải mẫu, đáp án do kho tự giải và kiểm bằng sympy."
%% ghi_chu: "Đề gốc ghi 'n là một số thực tùy ý' nhưng n được xác định theo m (n = 2m + 1); kho giữ nguyên chữ của đề."
\begin{ex}
Biết rằng hàm số $f(x)=\begin{cases}\dfrac{x^2+5x+6}{x+2}&\text{khi }x>-2\\mx+n&\text{khi }x\le-2\end{cases}$ liên tục trên $\mathbb R$ và $n$ là một số thực tùy ý. Khi đó
\choiceTF
{$\lim\limits_{x\to(-2)^+}f(x)=-2$.}
{Khi $m=1$ thì $n=2$.}
{\True Khi $m=2$ thì $n=5$.}
{\True Khi $m=3$ thì $n=7$.}
\loigiai{\sloppy Với $x>-2$: $f(x)=\dfrac{(x+2)(x+3)}{x+2}=x+3$, nên $\lim\limits_{x\to(-2)^+}f(x)=1$. Liên tục tại $-2$: $-2m+n=1\Rightarrow n=2m+1$.
a) Giới hạn bằng $1$, không phải $-2$. Sai.
b) $m=1\Rightarrow n=3$. Sai.
c) $m=2\Rightarrow n=5$. Đúng.
d) $m=3\Rightarrow n=7$. Đúng.}
\end{ex}
